\chapter{Security and Trust}
\label{ch:security}

Gert is designed to operate in regulated, high-trust environments where
operational actions must be auditable, accountable, and governed by explicit
policy.  Security is not an add-on; it is foundational to the system's
architecture.  This chapter defines the threat model, trust boundaries,
process isolation mechanisms, credential handling, transport security, RBAC
model, audit non-repudiation, and supply chain security.

% ─────────────────────────────────────────────────────────────────────────────
\section{Threat Model}
\label{sec:threat-model}
% ─────────────────────────────────────────────────────────────────────────────

Gert addresses the following threat categories, which are common in
operational runbook systems and have real-world precedent in systems like AWS
Systems Manager, PagerDuty Runbook Automation, and Rundeck~\cite{aws-ssm-runbooks,pagerduty-automation,rundeck-docs}.

\begin{table}[ht]
\centering
\begin{tabular}{p{0.35\linewidth}p{0.35\linewidth}p{0.20\linewidth}}
\hline
\textbf{Threat} & \textbf{Mitigation} & \textbf{Section Ref} \\
\hline
Runbook author executes unauthorized commands &
Command allowlist/denylist enforced at runtime core &
§\ref{sec:threat-model}, §11 \\
\hline
Malicious extension injects code into host process &
Out-of-process extension isolation; cryptographic signature verification &
§\ref{sec:extension-trust}, §4 \\
\hline
Tool or extension leaks credentials into trace &
Mandatory redaction for sensitive fields; \texttt{sensitive: true} input marking &
§\ref{sec:credential-handling} \\
\hline
Unauthorized user triggers runbook execution &
RBAC enforced at \texttt{gert serve} ingress &
§\ref{sec:rbac} \\
\hline
Man-in-the-middle attack on JSON-RPC transport &
TLS required for all non-localhost listeners &
§\ref{sec:transport-security} \\
\hline
Replay attack on approval decision &
Approval tokens signed with Ed25519; \texttt{approvalId} is single-use &
§\ref{sec:rbac}, §11 \\
\hline
Compromised tool binary executes untrusted code &
Tool binary checksum verification before invocation &
§\ref{sec:supply-chain} \\
\hline
\end{tabular}
\caption{Threat model and mitigations for gert.}
\label{tab:threat-model}
\end{table}

This threat model is derived from NIST SP 800-53 security controls~\cite{nist-sp800-53}
and ISO/IEC 27001 guidance on privileged access management and audit logging~\cite{iso27001}.

\begin{figure}[H]
\centering
\begin{tikzpicture}[scale=0.85, transform shape]
  % ── Pipeline stages ──────────────────────────────────────────────────
  \node[stepbox, minimum width=2.6cm, minimum height=0.8cm] (load)
    at (0,0) {Runbook\\Load};
  \node[stepbox, minimum width=2.6cm, minimum height=0.8cm,
    right=2.0cm of load] (runstart) {Run\\Start};
  \node[stepbox, minimum width=2.6cm, minimum height=0.8cm,
    right=2.0cm of runstart] (dispatch) {Step\\Dispatch};
  \node[stepbox, minimum width=2.6cm, minimum height=0.8cm,
    right=2.0cm of dispatch] (execute) {Execute\\Step};

  % ── Main flow arrows ─────────────────────────────────────────────────
  \draw[gertarrow] (load)     -- (runstart);
  \draw[gertarrow] (runstart) -- (dispatch);
  \draw[gertarrow] (dispatch) -- (execute);

  % ── Governance check labels above each stage ──────────────────────────
  \node[catbox, minimum width=2.6cm, above=0.9cm of load,
    font=\scriptsize\bfseries] (c1)
    {deny\_env\_vars\\ext signature\\schema valid};
  \node[catbox, minimum width=2.6cm, above=0.9cm of runstart,
    font=\scriptsize\bfseries] (c2)
    {requires\_approval\\actor role\\pre-flight policy};
  \node[catbox, minimum width=2.6cm, above=0.9cm of dispatch,
    font=\scriptsize\bfseries] (c3)
    {allowlist / denylist\\output redaction\\step-level rules};

  \draw[gertarrow, dashed, black!50] (c1) -- (load);
  \draw[gertarrow, dashed, black!50] (c2) -- (runstart);
  \draw[gertarrow, dashed, black!50] (c3) -- (dispatch);

  % ── BLOCKED exits below each stage ───────────────────────────────────
  \node[stepbox, fill=black!18, minimum width=2.6cm, below=1.0cm of load,
    font=\small\bfseries] (b1) {BLOCKED};
  \node[stepbox, fill=black!18, minimum width=2.6cm, below=1.0cm of runstart,
    font=\small\bfseries] (b2) {BLOCKED};
  \node[stepbox, fill=black!18, minimum width=2.6cm, below=1.0cm of dispatch,
    font=\small\bfseries] (b3) {BLOCKED};

  \draw[gertarrow, dashed] (load)     -- (b1)
    node[midway, right, align=center, font=\scriptsize] {policy\\fail};
  \draw[gertarrow, dashed] (runstart) -- (b2)
    node[midway, right, align=center, font=\scriptsize] {approval\\denied};
  \draw[gertarrow, dashed] (dispatch) -- (b3)
    node[midway, right, align=center, font=\scriptsize] {rule\\violation};
\end{tikzpicture}
\caption{Governance enforcement points across the gert execution pipeline.
  Each stage has dedicated policy checks; failure at any checkpoint halts
  execution and emits a \texttt{governance/blocked} trace event.}
\label{fig:governance-enforcement-points}
\end{figure}

% ─────────────────────────────────────────────────────────────────────────────
\section{Extension Trust Model}
\label{sec:extension-trust}
% ─────────────────────────────────────────────────────────────────────────────

Extensions operate under one of three trust levels, declared explicitly in the
extension manifest and evaluated by the host before any process is started.

\subsection{Trust Levels}

\paragraph{Local trusted.}
Installed by the machine owner (via package manager or explicit installation to
a system-wide extension directory).  These extensions run as child processes of
the current user, inherit the user's capability set (subject to manifest
restriction), and have access to all host capabilities declared in their
manifest.  No signature verification is performed; trust is conferred by the
installation act.

Example: an extension installed via \texttt{apt-get} or \texttt{brew} into
\texttt{/usr/local/lib/gert-extensions/}.

\paragraph{Project-scoped.}
Declared in the runbook project directory (e.g.,
\texttt{.gert/extensions.yaml}) or inline in the runbook's \texttt{extensions:}
field.  These extensions run in a restricted sandbox: only capabilities
explicitly granted by the workspace governance policy are forwarded.  File and
network access are constrained to declared paths and hosts.  Signature
verification is \emph{optional} but recommended.

Use case: a team-maintained extension for Acme-specific incident tooling,
distributed via the runbook repository.

\paragraph{Remote/signed.}
Extensions distributed via package repository or downloaded from the internet.
These extensions \emph{must} carry a cryptographic signature verified against a
trust store.  The signature covers the extension manifest and binary.  If
verification fails, the extension is rejected before any subprocess is started.

Signature format is defined in Section~\ref{sec:extension-signature}.

\subsection{Extension Signature Format}
\label{sec:extension-signature}

Remote/signed extensions must include a \texttt{signature:} block in
\texttt{gert-extension.yaml}:

\begin{minted}{yaml}
# gert-extension.yaml — signed extension
apiVersion: extension/v1

meta:
  name: acme.incident-pack
  version: "1.2.0"

signature:
  algorithm: ed25519
  publicKeyFingerprint: "sha256:abcd1234..."
  signatureData: "base64-encoded-signature"
  signedFields: ["meta", "capabilities", "entry-point", "transport"]

trust_chain:
  - fingerprint: "sha256:abcd1234..."
    issuer: "Acme Platform Team CA"
    notBefore: "2024-01-01T00:00:00Z"
    notAfter: "2025-01-01T00:00:00Z"
\end{minted}

\paragraph{Signature verification algorithm.}
The signature is computed over the canonical JSON representation of the fields
listed in \texttt{signedFields}.  The host:

\begin{enumerate}
  \item Extracts the public key fingerprint from the \texttt{signature} block.
  \item Looks up the corresponding public key in the trust store
        (\texttt{\$HOME/.gert/trusted-keys/} or
        \texttt{/etc/gert/trusted-keys/}).
  \item Reconstructs the canonical JSON of the signed fields (sorted keys,
        no whitespace).
  \item Verifies the Ed25519 signature over the canonical JSON bytes.
  \item Checks that the current time falls within the \texttt{notBefore} and
        \texttt{notAfter} window of the trust chain.
\end{enumerate}

If any step fails, the extension is rejected with a
\texttt{governance/extensionRejected} trace event and reason
\texttt{signature\_verification\_failed}.

\subsection{Extension Verification Command}

The \texttt{gert extension verify} command allows operators to manually verify
an extension before installation:

\begin{minted}{bash}
$ gert extension verify ./acme-extension/gert-extension.yaml \
    --trust-store /etc/gert/trusted-keys/
Extension: acme.incident-pack v1.2.0
Signature: valid (ed25519, expires 2025-01-01)
Issuer: Acme Platform Team CA
Capabilities: tool-registration, network
Trust level: remote/signed
Verification: PASSED
\end{minted}

% ─────────────────────────────────────────────────────────────────────────────
\section{Process Isolation}
\label{sec:process-isolation}
% ─────────────────────────────────────────────────────────────────────────────

Extensions and tools always run as separate OS processes.  They are never
loaded as shared libraries or executed in-process.  This isolation boundary
ensures that a crash, hang, or malicious extension cannot compromise the host
process.

\subsection{Isolation Mechanisms}

\paragraph{Separate OS process.}
As defined in §4, extensions are forked as child processes with their own
address space, file descriptor table, and signal handlers.

\paragraph{Filesystem access control.}
Extensions with \texttt{capability/file-read} or \texttt{capability/file-write}
declare the set of filesystem paths they may access via
\texttt{file-read-paths} and \texttt{file-write-paths} in the manifest.  The
host enforces these constraints:

\begin{itemize}
  \item On OpenBSD: via \texttt{unveil(2)} to restrict visible paths.
  \item On Linux: via seccomp-bpf to filter \texttt{open(2)} and
        \texttt{openat(2)} syscalls.
  \item On all platforms: by validating all path arguments in
        \texttt{file-read} and \texttt{file-write} RPC methods before dispatch.
\end{itemize}

Attempts to access undeclared paths return a capability-denied error.

\paragraph{Network access control.}
Extensions with \texttt{capability/network} declare the set of hostnames they
may connect to via \texttt{network-hosts}.  The host validates outbound
connections:

\begin{itemize}
  \item Before dispatching an extension RPC that specifies a target URL, the
        host checks that the hostname matches one of the declared
        \texttt{network-hosts} patterns.
  \item Extensions may not listen for inbound connections.
  \item Wildcard patterns are supported: \texttt{*.acme.example} permits
        \texttt{api.acme.example}, \texttt{prod.acme.example}, etc.
\end{itemize}

\paragraph{Environment variable isolation.}
Extensions receive ONLY the environment variables explicitly forwarded by the
runbook's \texttt{governance.allowed\_env} list.  Host environment variables
that match \texttt{governance.deny\_env\_vars} patterns are never forwarded,
even if they appear in \texttt{allowed\_env}.

The denylist takes precedence: if \texttt{AWS\_SECRET\_ACCESS\_KEY} is denied,
it is stripped from the extension's environment regardless of other policy.

\paragraph{Crash isolation.}
If an extension process crashes (exits with non-zero status or is killed by a
signal), the host:

\begin{enumerate}
  \item Emits an \texttt{extension.crashed} trace event with exit code and
        signal number.
  \item Marks the extension as unavailable for the remainder of the run.
  \item Cancels all in-flight invocations to that extension with error code
        \texttt{-32099} (\texttt{extension\_crashed}).
  \item Continues the run: the host process does not terminate.
\end{enumerate}

% ─────────────────────────────────────────────────────────────────────────────
\section{Credential and Secret Handling}
\label{sec:credential-handling}
% ─────────────────────────────────────────────────────────────────────────────

Gert does not store secrets.  It acts as a pass-through: secrets are
resolved from input providers at runtime, used for a single step invocation,
and then discarded.  Secrets never appear in trace logs or captured output.

\subsection{Secret Resolution and Redaction}

\paragraph{Input provider sensitivity marking.}
Input provider definitions declare which resolved values are sensitive via
\texttt{sensitive: true} in the provider schema:

\begin{minted}{yaml}
# vault.provider.yaml
outputs:
  credentials:
    type: object
    sensitive: true
\end{minted}

When the host resolves \texttt{from: vault.credentials}, it marks the returned
value as sensitive in memory.  This flag is respected by all downstream
consumers.

\paragraph{Tool sensitive inputs.}
Tool action definitions declare which input parameters accept sensitive data via
\texttt{sensitive\_inputs}:

\begin{minted}{yaml}
# my-tool.tool.yaml
actions:
  - name: deploy
    args:
      api_key:
        type: string
        required: true
    sensitive_inputs: [api_key]
\end{minted}

When the host prepares the tool invocation envelope, it marks the
\texttt{api\_key} field as sensitive.  The tool runtime redacts this field in
all trace events.

\paragraph{Mandatory redaction.}
All values marked \texttt{sensitive: true} are redacted before being written to
the trace or displayed in any adapter UI.  Redaction replaces the value with
the string \texttt{<redacted>}.

In addition, the runbook's \texttt{governance.redact} block defines regex
patterns applied to all captured output (stdout, stderr, JSON responses) before
storage.  See §11 for the full redaction specification.

\paragraph{Redacted \texttt{enum} member lists (C1).}
\texttt{enum} (design/gert/sections/03-schema-vnext.tex \S\ref{para:enum-constraint})
is \textbf{forbidden} on \texttt{type: secret} (\texttt{ENUM-001}): a secret's
value domain MUST NOT be enumerable, full stop. For every other declaration
that carries a sensitivity marking --- a tool \texttt{args:} entry with
\texttt{redact: true}, a name listed in \texttt{governance.sensitive\_inputs},
or a field matched by a \texttt{governance.redact} rule --- \texttt{enum} MAY
still be declared, but the following is normative and non-negotiable:

\begin{itemize}
  \item The member list \textbf{MUST NOT} appear in cleartext in any trace
        event, \texttt{--output=json} payload, \texttt{ValidatedPlan}
        record, or diagnostic. Every surface that would otherwise emit the
        member list emits the literal string \texttt{"<redacted>"} together
        with a \texttt{member\_count} integer instead.
  \item An \texttt{ENUM-008} (or \texttt{ENUM-009}) message for a redacted
        declaration \textbf{MUST NOT} enumerate permitted values and
        \textbf{MUST NOT} echo the rejected value; it MAY state only the
        declaration name and the fact that the value did not match one of
        \texttt{member\_count} permitted members. The rejected value itself
        is redacted per the ordinary rules above regardless of \texttt{enum}.
  \item This trades diagnostic quality on secret-adjacent fields for
        closing an oracle: an enum member list next to a redacted value is
        equivalent to publishing that value's entire domain, which defeats
        the redaction it sits beside.
\end{itemize}

\subsection{Secret Storage Policy}

Gert does NOT:

\begin{itemize}
  \item Store secrets in the runbook YAML.
  \item Write secrets to disk (trace files, state snapshots, etc.).
  \item Cache secrets across steps (unless explicitly declared by the input
        provider's \texttt{cache\_scope}).
  \item Transmit secrets over unencrypted transports.
\end{itemize}

Secrets must be resolved via input providers that integrate with secret
management systems (e.g., HashiCorp Vault, AWS Secrets Manager, 1Password).

% ─────────────────────────────────────────────────────────────────────────────
\section{Transport Security}
\label{sec:transport-security}
% ─────────────────────────────────────────────────────────────────────────────

Gert supports two primary transport modes for JSON-RPC communication:

\begin{enumerate}
  \item \textbf{stdio}: JSON-RPC over stdin/stdout of a child process.  No
        network exposure; inherently local.  No TLS required.
  \item \textbf{HTTP/WebSocket}: JSON-RPC over HTTP for \texttt{gert serve}
        mode.  Network exposure possible; TLS and authentication required for
        non-localhost.
\end{enumerate}

\subsection{TLS Requirements}

\paragraph{Localhost exception.}
If \texttt{gert serve} binds to \texttt{127.0.0.1} or \texttt{localhost}, TLS
is optional.  Connections are assumed to be local and trusted.

\paragraph{Non-localhost listeners.}
If \texttt{gert serve} binds to \texttt{0.0.0.0}, a public IP, or a hostname,
TLS is \emph{mandatory}.  The host will refuse to start without a valid TLS
certificate and private key.

Configuration:

\begin{minted}{bash}
$ gert serve --http 0.0.0.0:8443 \
    --tls-cert /etc/gert/tls/cert.pem \
    --tls-key /etc/gert/tls/key.pem
\end{minted}

The host performs standard TLS 1.3 handshake with the client.  TLS 1.2 is
supported for compatibility but TLS 1.1 and earlier are rejected.

\subsection{Client Authentication}

The host supports two authentication modes for \texttt{gert serve --http}:

\paragraph{Bearer token authentication.}
Clients include an \texttt{Authorization: Bearer <token>} header on all HTTP
requests.  The token is validated against a configured token store
(file-based, environment variable, or integration with an IdP).

Example:

\begin{minted}{bash}
$ gert serve --http :8443 --auth-mode bearer \
    --bearer-tokens /etc/gert/tokens.yaml
\end{minted}

The \texttt{tokens.yaml} file maps tokens to actor identities:

\begin{minted}{yaml}
# tokens.yaml
tokens:
  - token: "abc123..."
    actor: "ops@acme.example"
    roles: ["operator", "viewer"]
\end{minted}

\paragraph{Mutual TLS (mTLS).}
The host requires clients to present a valid X.509 client certificate signed by
a trusted CA.  The client certificate's subject CN or SAN is used as the actor
identity.

Configuration:

\begin{minted}{bash}
$ gert serve --http :8443 --auth-mode mtls \
    --tls-client-ca /etc/gert/tls/client-ca.pem
\end{minted}

mTLS is the recommended authentication mode for production deployments and
enterprise environments.

\subsection{WebSocket Security}

WebSocket connections for the event stream (§13) inherit the security context
of the HTTP server:

\begin{itemize}
  \item If the HTTP server uses TLS, WebSocket connections are encrypted
        (wss://).
  \item Client authentication applies: bearer tokens or client certificates are
        validated before the WebSocket upgrade.
  \item Once upgraded, the WebSocket connection is bound to the authenticated
        actor identity for the lifetime of the connection.
\end{itemize}

\subsection{Extension JSON-RPC Transport}

Extensions communicate with the host via stdio JSON-RPC (§4).  This transport
is not exposed to the network and does not require TLS.  The process boundary
is the trust boundary: the host trusts that OS-level process isolation prevents
interception of stdin/stdout pipes.

% ─────────────────────────────────────────────────────────────────────────────
\section{RBAC for \texttt{gert serve}}
\label{sec:rbac}
% ─────────────────────────────────────────────────────────────────────────────

When gert runs in server mode (\texttt{gert serve --http}), it enforces
role-based access control (RBAC) on all API operations.  Actors are assigned
roles, and each role grants a specific set of permissions.

\subsection{Actor Identity}

An \textbf{actor} is the authenticated identity making a request.  Actor
identity is derived from:

\begin{itemize}
  \item The bearer token (mapped to an actor ID via the token store), or
  \item The client certificate subject CN/SAN (in mTLS mode).
\end{itemize}

The actor identity is recorded in all trace events where a human or external
system initiates an action (run start, approval submission, cancellation).

\subsection{Roles}

Gert defines three standard roles:

\paragraph{Viewer.}
Read-only access.  Permissions:

\begin{itemize}
  \item List runs (\texttt{GET /runs})
  \item Get run details (\texttt{GET /runs/:id})
  \item Subscribe to event stream (\texttt{GET /ws} or \texttt{GET /events})
  \item Read trace files
\end{itemize}

Cannot start runs, approve, or cancel.

\paragraph{Operator.}
Execution permissions.  Includes all \textbf{viewer} permissions plus:

\begin{itemize}
  \item Start a run (\texttt{POST /exec/start})
  \item Cancel a run (\texttt{POST /exec/cancel})
  \item Submit approvals (\texttt{POST /approvals})
\end{itemize}

Cannot modify governance policy or manage roles.

\paragraph{Admin.}
Full control.  Includes all \textbf{operator} permissions plus:

\begin{itemize}
  \item Modify workspace governance policy
  \item Manage role assignments (add/remove actors from roles)
  \item Configure extension trust store
\end{itemize}

\subsection{Role Assignment}

Roles are assigned in the workspace config (\texttt{gert-project.yaml}):

\begin{minted}{yaml}
# gert-project.yaml
access:
  roles:
    - role: admin
      actors:
        - "alice@acme.example"
        - "bob@acme.example"
    - role: operator
      actors:
        - "ops-team@acme.example"
        - "oncall@acme.example"
    - role: viewer
      actors:
        - "*@acme.example"  # All Acme employees
\end{minted}

Actor patterns support wildcards for organization-wide grants.

\subsection{Approval Authorizer Role}

The \texttt{governance.require\_approval} field in the runbook declares the
role required to approve a step:

\begin{minted}{yaml}
steps:
  - id: delete_database
    cli: "kubectl delete database prod-db"
    approvals:
      min: 1
      roles: ["dba", "admin"]
\end{minted}

When the run reaches this step, it emits an
\texttt{approval/required} trace event.  The step is paused.  Only actors with
the \texttt{dba} or \texttt{admin} role may submit a valid approval.

The approval submission includes:

\begin{minted}{go}
POST /approvals
{
  "runId": "run-abc123",
  "stepId": "delete_database",
  "decision": "approve",
  "actorId": "alice@acme.example",
  "signature": "<ed25519-signature-of-approval-payload>",
  "timestamp": "2026-04-18T12:00:00Z"
}
\end{minted}

The host verifies:

\begin{enumerate}
  \item The actor is authenticated and has one of the required roles.
  \item The signature is valid (Ed25519 over the canonical JSON of
        \texttt{runId}, \texttt{stepId}, \texttt{decision}, \texttt{actorId},
        \texttt{timestamp}).
  \item The \texttt{timestamp} is within a 5-minute clock skew window.
  \item The \texttt{runId} and \texttt{stepId} match an active approval gate.
  \item This approval has not been submitted before (replay protection).
\end{enumerate}

If all checks pass, the approval is recorded in the trace
(\texttt{approval/received}) and the step proceeds if the minimum approval
count is reached.

% ─────────────────────────────────────────────────────────────────────────────
\section{Audit Non-Repudiation}
\label{sec:audit-nonrepudiation}
% ─────────────────────────────────────────────────────────────────────────────

Non-repudiation ensures that governance decisions — especially approvals and
rejections — cannot be denied by the actor who made them.  This is critical for
regulatory compliance (SOC2, ISO27001, HIPAA)~\cite{soc2-aicpa,iso27001,hipaa-security-rule}.

\subsection{Governance Decision Recording}

Every governance decision event (\texttt{governance/approval\_received},
\texttt{governance/approval\_rejected}, \texttt{governance/commandBlocked})
must record:

\begin{itemize}
  \item Actor identity (authenticated user or service principal)
  \item Timestamp (RFC 3339 with timezone)~\cite{rfc3339}
  \item Decision payload (what was approved/rejected)
  \item Signature (Ed25519 over the canonical decision payload)
\end{itemize}

Example trace event:

\begin{minted}{json}
{
  "eventType": "governance/approval_received",
  "timestamp": "2026-04-18T12:00:00Z",
  "runId": "run-abc123",
  "stepId": "delete_database",
  "actor": "alice@acme.example",
  "decision": "approve",
  "signature": "ed25519:abcd1234...",
  "signatureAlgorithm": "ed25519",
  "signaturePublicKey": "sha256:pubkey-fingerprint"
}
\end{minted}

The signature covers the fields \texttt{runId}, \texttt{stepId},
\texttt{actor}, \texttt{decision}, and \texttt{timestamp}.  The public key is
looked up from the actor's identity record or client certificate.

\subsection{Trace File Integrity (HMAC Chaining)}

The JSONL trace file (§12) is append-only, but this alone does not prevent
tampering.  To provide cryptographic integrity, the host optionally computes an
HMAC chain over the trace:

\begin{enumerate}
  \item At run start, initialize \texttt{H\_0 = HMAC-SHA256(key, runId)}.
  \item After writing each trace event \texttt{E\_i}, compute
        \texttt{H\_i = HMAC-SHA256(key, H\_\{i-1\} || E\_i)}.
  \item Store \texttt{H\_i} in the event as \texttt{integrityHash}.
  \item At run completion, emit a final \texttt{run.completed} event with
        \texttt{finalHash: H\_N}.
\end{enumerate}

Verification: a verifier recomputes the HMAC chain from \texttt{H\_0} forward.
If any event is modified or deleted, the chain breaks.

The HMAC key is derived from the workspace master key or provided via
\texttt{--integrity-key}.

\subsection{SHA256 Evidence Capture}

As defined in §12, gert captures SHA256 hashes of all evidence artifacts
(attachments, screenshots, CLI output files).  These hashes are written to the
trace event immediately after the artifact is captured, providing a
cryptographic anchor for the artifact's integrity.

The combination of HMAC trace chaining + SHA256 evidence hashes provides
end-to-end non-repudiation for the entire run.

% ─────────────────────────────────────────────────────────────────────────────
\section{Supply Chain Security}
\label{sec:supply-chain}
% ─────────────────────────────────────────────────────────────────────────────

Tools and extensions are executable code distributed via package managers,
repositories, or direct download.  Supply chain attacks (e.g., compromised
binaries, backdoored dependencies) are a real threat in operational
environments.  Gert provides checksum verification and signature validation
to mitigate these risks.

\subsection{Tool Binary Checksums}

Tool definitions may declare a \texttt{checksum:} field in
\texttt{.tool.yaml}:

\begin{minted}{yaml}
# kubectl.tool.yaml
meta:
  name: kubectl
  binary: kubectl

checksum:
  algorithm: sha256
  digest: "a3b4c5d6e7f8..."
\end{minted}

Before invoking the tool, the host:

\begin{enumerate}
  \item Resolves the binary path (via \texttt{PATH} lookup or absolute path).
  \item Computes the SHA256 hash of the binary file.
  \item Compares the computed hash to \texttt{checksum.digest}.
  \item If the hashes do not match, the tool invocation is blocked with a
        \texttt{governance/toolChecksumMismatch} trace event.
\end{enumerate}

\subsection{Tool Package Digests and Lock File}
\label{sec:supply-chain-package-digests}

In addition to the per-binary checksum above, tool \textbf{packages}
(directories bundling \texttt{.tool.yaml} definitions behind
\texttt{gert-package.yaml}, \texttt{design/gert/sections/06-tool-runtime.tex}
\S Tool Packages) are integrity-protected at the package level:

\begin{itemize}
  \item Every package has a deterministic, order-free \textbf{package
        digest} computed over its exported API surface (manifest, exported
        tool files, and everything transitively reachable via
        \texttt{impl.path} and package-internal includes) --- see
        \S Digest, Evidence, Trace, Resume in the tool runtime chapter for
        the exact construction.
  \item The full package set for a run has a single \textbf{catalog digest}
        that proves two runs saw the same tool universe.
  \item \texttt{.gert/packages.lock.json}
        (\texttt{design/gert/schemas/package-lock.v1.schema.json},
        \texttt{apiVersion: package-lock/v1}) is a generated, committed,
        frozen integrity record of both --- not a dependency solver output.
        Lock presence is optional for \texttt{gert run}; when present, a
        mismatch is \texttt{PKG-009}.
  \item On \texttt{gert exec --resume}, package/catalog digest mismatch is a
        hard refusal (\texttt{PKG-009}), overridable only via
        \texttt{--allow-package-drift}, which emits an auditable
        \texttt{governance/packageDriftAccepted} trace event --- see
        \texttt{design/gert/sections/13-evidence-tracing-resumption.tex}
        \S Package Resumption Contract.
\end{itemize}

\subsection{\texttt{gert verify} Command}

The \texttt{gert verify} command validates all tools, packages, and
extensions in the workspace before execution:

\begin{minted}{bash}
$ gert verify --workspace .
Verifying tools...
  ✓ kubectl: checksum verified (sha256:a3b4c5...)
  ✓ curl: no checksum declared (skipped)
  ✗ jq: checksum mismatch (expected sha256:abc123, got sha256:def456)

Verifying packages...
  ✓ acme.incident-tools@1.4.2: package digest matches lock (sha256:9f2e1a...)
  ⚠ acme.audit-pack@0.3.0: external root (vendor path resolves outside workspace)
  ✗ acme.legacy-tools@2.0.0: package digest mismatch (lock sha256:aa11.., disk sha256:bb22..)

Verifying extensions...
  ✓ acme.incident-pack: signature valid (ed25519, expires 2025-01-01)
  ✗ debug-tools: no signature (required for remote extensions)

Verification: FAILED (2 errors, 1 warning)
\end{minted}

This command should be run in CI/CD pipelines before deploying runbooks to
production.

\subsubsection{Package verification}
\label{sec:supply-chain-packages}

Per Barbara's ruling
\texttt{.squad/decisions/archive/barbara-tool-packages-architecture-ruling.md}
(AR-TP-6, AR-TP-7, ratified 2026-08-09), \texttt{gert verify} additionally:

\begin{enumerate}
  \item Recomputes every package digest and the catalog digest from the
        current filesystem, using the raw-byte SHA-256 construction in
        \texttt{design/gert/sections/06-tool-runtime.tex} \S Digest,
        Evidence, Trace, Resume.
  \item If \texttt{.gert/packages.lock.json} is present, compares every
        recomputed digest against the corresponding \texttt{packages[].digest},
        \texttt{packages[].exports[].digest}, and \texttt{catalogDigest}
        entries. Any mismatch is reported as a distinct \texttt{PKG-009}
        finding; the lock file itself is validated against
        \texttt{design/gert/schemas/package-lock.v1.schema.json}.
  \item Reports every package whose root resolves outside the workspace
        (\S Secure Path Resolution, \texttt{PKG-W003}) as a distinct
        ``external root'' finding --- a warning, not a failure, but one that
        \textbf{MUST} be surfaced so operators can audit what is being
        loaded from outside the workspace tree.
  \item Applies the same symlink/junction/reparse-point resolution and
        containment rules used at plan time (\S Secure Path Resolution
        \S Symlinks, junctions, reparse points); an escaping link is
        \texttt{PKG-008}.
\end{enumerate}

\subsection{Extension Signature Verification}

As defined in §\ref{sec:extension-signature}, remote extensions must carry a
cryptographic signature.  The signature verification process ensures that:

\begin{itemize}
  \item The extension binary and manifest have not been tampered with.
  \item The extension was signed by a trusted authority (public key in the
        trust store).
  \item The signature has not expired (\texttt{notAfter} check).
\end{itemize}

If verification fails, the extension is rejected before any subprocess is
started.

% ─────────────────────────────────────────────────────────────────────────────
\section{Security Recommendations}
% ─────────────────────────────────────────────────────────────────────────────

For production deployments, the following security practices are recommended:

\begin{enumerate}
  \item \textbf{Enable TLS for all \texttt{gert serve} deployments} — even
        internal networks can be compromised.
  \item \textbf{Use mTLS for client authentication} — bearer tokens are
        vulnerable to theft; client certificates provide stronger identity.
  \item \textbf{Enable HMAC trace chaining} — provides cryptographic integrity
        for audit trails.
  \item \textbf{Declare tool binary checksums} — prevents execution of
        tampered binaries.
  \item \textbf{Require signatures for all remote extensions} — do not trust
        unsigned code from the internet.
  \item \textbf{Run \texttt{gert verify} in CI} — catch integrity violations
        before production.
  \item \textbf{Use least-privilege capability grants} — extensions should
        request only the capabilities they need; admins should audit and deny
        excessive requests.
  \item \textbf{Regularly rotate approval keys} — the Ed25519 keys used for
        signing approvals should be rotated every 90 days.
\end{enumerate}

These recommendations align with NIST SP 800-53 controls for privileged access
management (AC-2, AC-6), audit and accountability (AU-2, AU-9), and system
integrity (SI-7)~\cite{nist-sp800-53}.
